% Generated by Sphinx.
\documentclass[a4paper,10pt,english]{manual}
\usepackage[utf8]{inputenc}
\usepackage[T1]{fontenc}
\usepackage{babel}
\usepackage{times}
\usepackage[Bjarne]{fncychap}
\usepackage{longtable}
\usepackage{sphinx}


\title{Workshop - FOSS4G routing with pgRouting tools, OpenStreetMap road data and GeoExt Manual}
\date{July 05, 2010}
\release{1}
\author{Daniel Kastl, Frédéric Junod}
\newcommand{\sphinxlogo}{}
\renewcommand{\releasename}{Release}
\makeindex


\makeatletter
\def\PYG@reset{\let\PYG@it=\relax \let\PYG@bf=\relax%
    \let\PYG@ul=\relax \let\PYG@tc=\relax%
    \let\PYG@bc=\relax \let\PYG@ff=\relax}
\def\PYG@tok#1{\csname PYG@tok@#1\endcsname}
\def\PYG@toks#1+{\ifx\relax#1\empty\else%
    \PYG@tok{#1}\expandafter\PYG@toks\fi}
\def\PYG@do#1{\PYG@bc{\PYG@tc{\PYG@ul{%
    \PYG@it{\PYG@bf{\PYG@ff{#1}}}}}}}
\def\PYG#1#2{\PYG@reset\PYG@toks#1+\relax+\PYG@do{#2}}

\def\PYG@tok@gd{\def\PYG@tc##1{\textcolor[rgb]{0.63,0.00,0.00}{##1}}}
\def\PYG@tok@gu{\let\PYG@bf=\textbf\def\PYG@tc##1{\textcolor[rgb]{0.50,0.00,0.50}{##1}}}
\def\PYG@tok@gt{\def\PYG@tc##1{\textcolor[rgb]{0.00,0.25,0.82}{##1}}}
\def\PYG@tok@gs{\let\PYG@bf=\textbf}
\def\PYG@tok@gr{\def\PYG@tc##1{\textcolor[rgb]{1.00,0.00,0.00}{##1}}}
\def\PYG@tok@cm{\let\PYG@it=\textit\def\PYG@tc##1{\textcolor[rgb]{0.25,0.50,0.56}{##1}}}
\def\PYG@tok@vg{\def\PYG@tc##1{\textcolor[rgb]{0.73,0.38,0.84}{##1}}}
\def\PYG@tok@m{\def\PYG@tc##1{\textcolor[rgb]{0.13,0.50,0.31}{##1}}}
\def\PYG@tok@mh{\def\PYG@tc##1{\textcolor[rgb]{0.13,0.50,0.31}{##1}}}
\def\PYG@tok@cs{\def\PYG@tc##1{\textcolor[rgb]{0.25,0.50,0.56}{##1}}\def\PYG@bc##1{\colorbox[rgb]{1.00,0.94,0.94}{##1}}}
\def\PYG@tok@ge{\let\PYG@it=\textit}
\def\PYG@tok@vc{\def\PYG@tc##1{\textcolor[rgb]{0.73,0.38,0.84}{##1}}}
\def\PYG@tok@il{\def\PYG@tc##1{\textcolor[rgb]{0.13,0.50,0.31}{##1}}}
\def\PYG@tok@go{\def\PYG@tc##1{\textcolor[rgb]{0.19,0.19,0.19}{##1}}}
\def\PYG@tok@cp{\def\PYG@tc##1{\textcolor[rgb]{0.00,0.44,0.13}{##1}}}
\def\PYG@tok@gi{\def\PYG@tc##1{\textcolor[rgb]{0.00,0.63,0.00}{##1}}}
\def\PYG@tok@gh{\let\PYG@bf=\textbf\def\PYG@tc##1{\textcolor[rgb]{0.00,0.00,0.50}{##1}}}
\def\PYG@tok@ni{\let\PYG@bf=\textbf\def\PYG@tc##1{\textcolor[rgb]{0.84,0.33,0.22}{##1}}}
\def\PYG@tok@nl{\let\PYG@bf=\textbf\def\PYG@tc##1{\textcolor[rgb]{0.00,0.13,0.44}{##1}}}
\def\PYG@tok@nn{\let\PYG@bf=\textbf\def\PYG@tc##1{\textcolor[rgb]{0.05,0.52,0.71}{##1}}}
\def\PYG@tok@no{\def\PYG@tc##1{\textcolor[rgb]{0.38,0.68,0.84}{##1}}}
\def\PYG@tok@na{\def\PYG@tc##1{\textcolor[rgb]{0.25,0.44,0.63}{##1}}}
\def\PYG@tok@nb{\def\PYG@tc##1{\textcolor[rgb]{0.00,0.44,0.13}{##1}}}
\def\PYG@tok@nc{\let\PYG@bf=\textbf\def\PYG@tc##1{\textcolor[rgb]{0.05,0.52,0.71}{##1}}}
\def\PYG@tok@nd{\let\PYG@bf=\textbf\def\PYG@tc##1{\textcolor[rgb]{0.33,0.33,0.33}{##1}}}
\def\PYG@tok@ne{\def\PYG@tc##1{\textcolor[rgb]{0.00,0.44,0.13}{##1}}}
\def\PYG@tok@nf{\def\PYG@tc##1{\textcolor[rgb]{0.02,0.16,0.49}{##1}}}
\def\PYG@tok@si{\let\PYG@it=\textit\def\PYG@tc##1{\textcolor[rgb]{0.44,0.63,0.82}{##1}}}
\def\PYG@tok@s2{\def\PYG@tc##1{\textcolor[rgb]{0.25,0.44,0.63}{##1}}}
\def\PYG@tok@vi{\def\PYG@tc##1{\textcolor[rgb]{0.73,0.38,0.84}{##1}}}
\def\PYG@tok@nt{\let\PYG@bf=\textbf\def\PYG@tc##1{\textcolor[rgb]{0.02,0.16,0.45}{##1}}}
\def\PYG@tok@nv{\def\PYG@tc##1{\textcolor[rgb]{0.73,0.38,0.84}{##1}}}
\def\PYG@tok@s1{\def\PYG@tc##1{\textcolor[rgb]{0.25,0.44,0.63}{##1}}}
\def\PYG@tok@gp{\let\PYG@bf=\textbf\def\PYG@tc##1{\textcolor[rgb]{0.78,0.36,0.04}{##1}}}
\def\PYG@tok@sh{\def\PYG@tc##1{\textcolor[rgb]{0.25,0.44,0.63}{##1}}}
\def\PYG@tok@ow{\let\PYG@bf=\textbf\def\PYG@tc##1{\textcolor[rgb]{0.00,0.44,0.13}{##1}}}
\def\PYG@tok@sx{\def\PYG@tc##1{\textcolor[rgb]{0.78,0.36,0.04}{##1}}}
\def\PYG@tok@bp{\def\PYG@tc##1{\textcolor[rgb]{0.00,0.44,0.13}{##1}}}
\def\PYG@tok@c1{\let\PYG@it=\textit\def\PYG@tc##1{\textcolor[rgb]{0.25,0.50,0.56}{##1}}}
\def\PYG@tok@kc{\let\PYG@bf=\textbf\def\PYG@tc##1{\textcolor[rgb]{0.00,0.44,0.13}{##1}}}
\def\PYG@tok@c{\let\PYG@it=\textit\def\PYG@tc##1{\textcolor[rgb]{0.25,0.50,0.56}{##1}}}
\def\PYG@tok@mf{\def\PYG@tc##1{\textcolor[rgb]{0.13,0.50,0.31}{##1}}}
\def\PYG@tok@err{\def\PYG@bc##1{\fcolorbox[rgb]{1.00,0.00,0.00}{1,1,1}{##1}}}
\def\PYG@tok@kd{\let\PYG@bf=\textbf\def\PYG@tc##1{\textcolor[rgb]{0.00,0.44,0.13}{##1}}}
\def\PYG@tok@ss{\def\PYG@tc##1{\textcolor[rgb]{0.32,0.47,0.09}{##1}}}
\def\PYG@tok@sr{\def\PYG@tc##1{\textcolor[rgb]{0.14,0.33,0.53}{##1}}}
\def\PYG@tok@mo{\def\PYG@tc##1{\textcolor[rgb]{0.13,0.50,0.31}{##1}}}
\def\PYG@tok@mi{\def\PYG@tc##1{\textcolor[rgb]{0.13,0.50,0.31}{##1}}}
\def\PYG@tok@kn{\let\PYG@bf=\textbf\def\PYG@tc##1{\textcolor[rgb]{0.00,0.44,0.13}{##1}}}
\def\PYG@tok@o{\def\PYG@tc##1{\textcolor[rgb]{0.40,0.40,0.40}{##1}}}
\def\PYG@tok@kr{\let\PYG@bf=\textbf\def\PYG@tc##1{\textcolor[rgb]{0.00,0.44,0.13}{##1}}}
\def\PYG@tok@s{\def\PYG@tc##1{\textcolor[rgb]{0.25,0.44,0.63}{##1}}}
\def\PYG@tok@kp{\def\PYG@tc##1{\textcolor[rgb]{0.00,0.44,0.13}{##1}}}
\def\PYG@tok@w{\def\PYG@tc##1{\textcolor[rgb]{0.73,0.73,0.73}{##1}}}
\def\PYG@tok@kt{\def\PYG@tc##1{\textcolor[rgb]{0.56,0.13,0.00}{##1}}}
\def\PYG@tok@sc{\def\PYG@tc##1{\textcolor[rgb]{0.25,0.44,0.63}{##1}}}
\def\PYG@tok@sb{\def\PYG@tc##1{\textcolor[rgb]{0.25,0.44,0.63}{##1}}}
\def\PYG@tok@k{\let\PYG@bf=\textbf\def\PYG@tc##1{\textcolor[rgb]{0.00,0.44,0.13}{##1}}}
\def\PYG@tok@se{\let\PYG@bf=\textbf\def\PYG@tc##1{\textcolor[rgb]{0.25,0.44,0.63}{##1}}}
\def\PYG@tok@sd{\let\PYG@it=\textit\def\PYG@tc##1{\textcolor[rgb]{0.25,0.44,0.63}{##1}}}

\def\PYGZbs{\char`\\}
\def\PYGZus{\char`\_}
\def\PYGZob{\char`\{}
\def\PYGZcb{\char`\}}
\def\PYGZca{\char`\^}
% for compatibility with earlier versions
\def\PYGZat{@}
\def\PYGZlb{[}
\def\PYGZrb{]}
\makeatother

\begin{document}

\maketitle
\tableofcontents
\hypertarget{--doc-latex}{}


\resetcurrentobjects
\hypertarget{--doc-chapters/introduction}{}

\chapter{Introduction}
\paragraph{Abstract}

\href{http://www.pgrouting.org}{pgRouting} adds routing functionality to \href{http://www.postgis.org}{PostGIS}. This introductory workshop will show you how. It gives a practical example of how to use pgRouting with \href{http://www.openstreetmap.org}{OpenStreetMap} road network data. It explains the steps to prepare the data, make routing queries, assign costs and use \href{http://www.geoext.org}{GeoExt} to show your route in a web-mapping application.

Navigation for road networks requires complex routing algorithms that support turn restrictions and even time-dependent attributes. pgRouting is an extendible open-source library that provides a variety of tools for shortest path search as extension of PostgreSQL and PostGIS. The workshop will explain about shortest path search with pgRouting in real road networks and how the data structure is important to get faster results. Also you will learn about difficulties and limitations of pgRouting in GIS applications.

To give a practical example the workshop makes use of OpenStreetMap data of Barcelona. You will learn how to convert the data into the required format and how to calibrate the data with ``cost'' attributes. Furthermore we will explain the difference of the main routing algorithms ``Dijkstra'', ``A-Star'' and ``Shooting-Star''. By the end of the workshop you will have a good understanding of how to use pgRouting and how to get your network data prepared.

To learn how to get the output from rows and columns to be drawn on a map, we will build a basic map GUI with GeoExt. We listened to the students feedback of the last year's and want to guide you through the basic steps to build a simple browser application. Our goal is to make this as easy as possible, and to show that it's not difficult to integrate with other FOSS4G tools. For that reason we selected GeoExt, which is a JavaScript library providing the groundwork for creating web-mapping applications based on OpenLayers and Ext.
\paragraph{Presenter}
\begin{itemize}
\item {} 
\emph{Daniel Kastl} is founder and CEO of \href{http://georepublic.de}{Georepublic UG} and works in Germany and Japan. He is moderating and promoting the pgRouting community and development since 4 years, and he's an active OSM contributor in Japan.

\item {} 
\emph{Frédéric Junod} works at the Swiss office of \href{http://www.camptocamp.com}{Camptocamp} for about five years. He's an active developer of many open source GIS projects from the browser (GeoExt, OpenLayers) to the server world (MapFish, Shapely, TileCache).

\end{itemize}

Daniel and Frédéric are the auhtors of the previous pgRouting workshops, that have been held at FOSS4G events in Canada and South Africa and at local conferences in Japan.

\begin{notice}{note}{Note:}\begin{itemize}
\item {} 
Workshop level: intermediate

\item {} 
Attendee's previous knowledge: SQL (PostgreSQL, PostGIS), Javascript, HTML

\item {} 
Equipments: This workshops will make use of the GIS LiveDVD if possible. Otherwise it will require VirtualBox installed to load a virtual machine image.

\end{itemize}
\end{notice}
\paragraph{License}

This work is licensed under a \href{http://creativecommons.org/licenses/by-sa/3.0/}{Creative Commons Attribution-Share Alike 3.0 License}.

\includegraphics{license.png}

\resetcurrentobjects
\hypertarget{--doc-chapters/about}{}

\chapter{About}
\hypertarget{todo-27}{}
\begin{notice}{note}{Todo}

introduction paragraph for about
\end{notice}


\section{pgRouting}

pgRouting is an extension of PostGIS and adds routing functionality to PostGIS/PostgreSQL. pgRouting is a further development of pgDijkstra (by Camptocamp). It is currently developed and maintained by Orkney, JAPAN.

pgRouting provides functions for:
\begin{itemize}
\item {} 
Shortest Path Dikstra: routing algorithm without heuristics

\item {} 
Shortest Path A-Star: routing for large datasets (with heuristics)

\item {} 
Shortest Path Shooting-Star: routing with turn restrictions (with heuristics)

\item {} 
Traveling Salesperson Problem (TSP)

\item {} 
Driving Distance calculation (Isolines)

\end{itemize}

Advantages of the database routing approach are:
\begin{itemize}
\item {} 
Accessible by multiple clients through JDBC, ODBC, or directly using Pl/pgSQL. The clients can either be PCs or mobile devices.

\item {} 
Uses PostGIS for its geographic data format, which in turn uses OGC's data format Well Konwn Text (WKT) and Well Known Binary (WKB). This allows usage of existing open * data converters.

\item {} 
Open Source software like qGIS and uDig can modify the data/attributes,

\item {} 
Data changes can be reflected instantaneously through the routing engine. There is no need for precalculation.

\item {} 
The ``cost'' parameter can be dynamically calculated through SQL and its value can come from multiple fields or tables.

\end{itemize}

pgRouting project website: \href{http://www.pgrouting.org}{http://www.pgrouting.org}


\section{OpenStreetMap}

\emph{``OpenStreetMap is a project aimed squarely at creating and providing free geographic data such as street maps to anyone who wants them.'' ``The project was started because most maps you think of as free actually have legal or technical restrictions on their use, holding back people from using them in creative, productive or unexpected ways.''} \href{http://wiki.openstreetmap.org/index.php/Press}{http://wiki.openstreetmap.org/index.php/Press}

OpenStreetMap uses a topological data structure:
\begin{itemize}
\item {} 
Nodes are points with a geographic position.

\item {} 
Ways are lists of nodes, representing a polyline or polygon.

\item {} 
Relations are groups of nodes, ways and other relations which can be assigned certain properties.

\item {} 
Tags can be applied to nodes, ways or relations and consist of name=value pairs.

\item {} 
This is how nodes, ways and relations are described in the OpenStreetMap XML file:

\end{itemize}

\begin{Verbatim}[commandchars=@\[\]]
@textless[]?xml version='1.0' encoding='UTF-8'?@textgreater[]
@textless[]osm version='0.6' generator='xapi: OSM Extended API 2.0' ... @textgreater[]
  ...	
  @textless[]node id='255405560' lat='41.4917468' lon='2.0257695' version='1' 
  		changeset='19117' user='efrainlarrea' uid='32823' visible='true' 
  		timestamp='2008-04-02T17:40:07Z'@textgreater[]
  @textless[]/node@textgreater[]
  @textless[]node id='255405551' lat='41.4866740' lon='2.0302842' version='3' 
  		changeset='248452' user='efrainlarrea' uid='32823' visible='true' 
  		timestamp='2008-04-24T15:56:08Z'@textgreater[]
  @textless[]/node@textgreater[]
  @textless[]node id='255405552' lat='41.4868540' lon='2.0297863' version='1' 
  		changeset='19117' user='efrainlarrea' uid='32823' visible='true' 
  		timestamp='2008-04-02T17:40:07Z'@textgreater[]
  @textless[]/node@textgreater[]
  ...
  @textless[]way id='35419222' visible='true' timestamp='2009-06-03T21:49:11Z' 
  		version='1' changeset='1416898' user='Yodeima' uid='115931'@textgreater[]
    @textless[]nd ref='415466914'/@textgreater[]
    @textless[]nd ref='415466915'/@textgreater[]
    @textless[]tag k='highway' v='unclassified'/@textgreater[]
    @textless[]tag k='lanes' v='1'/@textgreater[]
    @textless[]tag k='name' v='Carrer del Progrés'/@textgreater[]
    @textless[]tag k='oneway' v='no'/@textgreater[]
  @textless[]/way@textgreater[]
  @textless[]way id='35419227' visible='true' timestamp='2009-06-14T20:37:55Z' 
  		version='2' changeset='1518775' user='Yodeima' uid='115931'@textgreater[]
    @textless[]nd ref='415472085'/@textgreater[]
    @textless[]nd ref='415472086'/@textgreater[]
    @textless[]nd ref='415472087'/@textgreater[]
    @textless[]tag k='highway' v='unclassified'/@textgreater[]
    @textless[]tag k='lanes' v='1'/@textgreater[]
    @textless[]tag k='name' v='carrer de la mecanica'/@textgreater[]
    @textless[]tag k='oneway' v='no'/@textgreater[]
  @textless[]/way@textgreater[]
  ...
  @textless[]relation id='903432' visible='true' timestamp='2010-05-06T08:36:54Z' 
  		version='1' changeset='4619553' user='ivansanchez' uid='5265'@textgreater[]
    @textless[]member type='way' ref='56426179' role='outer'/@textgreater[]
    @textless[]member type='way' ref='56426173' role='inner'/@textgreater[]
    @textless[]tag k='layer' v='0'/@textgreater[]
    @textless[]tag k='leisure' v='common'/@textgreater[]
    @textless[]tag k='name' v='Plaça Can Suris'/@textgreater[]
    @textless[]tag k='source' v='WMS shagrat.icc.cat'/@textgreater[]
    @textless[]tag k='type' v='multipolygon'/@textgreater[]
  @textless[]/relation@textgreater[]
  ...
@textless[]/osm@textgreater[]
\end{Verbatim}

The OSM data can be downloaded from OpenStreetMap website using an API (see  \href{http://wiki.openstreetmap.org/index.php/OSM\_Protocol\_Version\_0.5}{http://wiki.openstreetmap.org/index.php/OSM\_Protocol\_Version\_0.5}), or with some other OSM tools, for example JOSM editor.

Update: CloudMade offers extracts of maps from different places around the world. For South Africa go to  \href{http://download.cloudmade.com/africa/south\_africa}{http://download.cloudmade.com/africa/south\_africa}

Note: The API has a download size limitation, which can make it a bit inconvenient to download extensive areas with many features.


\section{osm2pgrouting}

When using the osm2pgrouting converter (see later), we take only nodes and ways of types and classes specified in ``mapconfig.xml'' file to be converted to pgRouting table format:

\begin{Verbatim}[commandchars=\\\{\}]
\PYG{c+cp}{\textless{}?xml version="1.0" encoding="UTF-8"?\textgreater{}}
\PYG{n+nt}{\textless{}configuration}\PYG{n+nt}{\textgreater{}}
  \PYG{n+nt}{\textless{}type} \PYG{n+na}{name=}\PYG{l+s}{"highway"} \PYG{n+na}{id=}\PYG{l+s}{"1"}\PYG{n+nt}{\textgreater{}}
    \PYG{n+nt}{\textless{}class} \PYG{n+na}{name=}\PYG{l+s}{"motorway"} \PYG{n+na}{id=}\PYG{l+s}{"101"} \PYG{n+nt}{/\textgreater{}}
    \PYG{n+nt}{\textless{}class} \PYG{n+na}{name=}\PYG{l+s}{"motorway\PYGZus{}link"} \PYG{n+na}{id=}\PYG{l+s}{"102"} \PYG{n+nt}{/\textgreater{}}
    \PYG{n+nt}{\textless{}class} \PYG{n+na}{name=}\PYG{l+s}{"motorway\PYGZus{}junction"} \PYG{n+na}{id=}\PYG{l+s}{"103"} \PYG{n+nt}{/\textgreater{}}
    ...
    \PYG{n+nt}{\textless{}class} \PYG{n+na}{name=}\PYG{l+s}{"road"} \PYG{n+na}{id=}\PYG{l+s}{"100"} \PYG{n+nt}{/\textgreater{}}
  \PYG{n+nt}{\textless{}/type\textgreater{}}    
  \PYG{n+nt}{\textless{}type} \PYG{n+na}{name=}\PYG{l+s}{"junction"} \PYG{n+na}{id=}\PYG{l+s}{"4"}\PYG{n+nt}{\textgreater{}}
    \PYG{n+nt}{\textless{}class} \PYG{n+na}{name=}\PYG{l+s}{"roundabout"} \PYG{n+na}{id=}\PYG{l+s}{"401"} \PYG{n+nt}{/\textgreater{}}
  \PYG{n+nt}{\textless{}/type\textgreater{}}  
\PYG{n+nt}{\textless{}/configuration\textgreater{}}
\end{Verbatim}

Detailed description of all possible types and classes can be found here:  \href{http://wiki.openstreetmap.org/index.php/Map\_features}{http://wiki.openstreetmap.org/index.php/Map\_features}.

For Cape Town the OpenStreetMap data is very comprehensive with many details. A compilation of the greater Cape Town area created with JOSM is available as capetown\_20080829.osm.


\section{GeoExt}
\hypertarget{todo-28}{}
\begin{notice}{note}{Todo}

GeoExt paragraph
\end{notice}

\resetcurrentobjects
\hypertarget{--doc-chapters/installation}{}

\chapter{Installation and Requirements}

\resetcurrentobjects
\hypertarget{--doc-chapters/osm2pgrouting}{}

\chapter{osm2pgrouting Import Tool}

This tool makes it easy to import OpenStreetMap data and use it with pgRouting. It creates topology automatically and creates tables for feature types and road classes. osm2pgrouting was primarily written by Daniel Wendt and is now hosted on the pgRouting project site: \href{http://pgrouting.postlbs.org/wiki/tools/osm2pgrouting}{http://pgrouting.postlbs.org/wiki/tools/osm2pgrouting}

How to install (Ubuntu 8.04)

Check out the latest version from SVN repository:

svn checkout \href{http://pgrouting.postlbs.org/svn/pgrouting/tools/osm2pgrouting/trunk}{http://pgrouting.postlbs.org/svn/pgrouting/tools/osm2pgrouting/trunk} osm2pgrouting
Required packages/libraries:
\begin{enumerate}
\item {} 
PostgreSQL 2. PostGIS 3. pgRouting 4. Boost library 5. Expat library 6. libpq library

\end{enumerate}

Note: if you already compiled pgRouting point 1. to 4. should already be installed.

Then compile

cd osm2pgrouting
make
How to use
\begin{enumerate}
\item {} 
First you need to create a database and add PostGIS and pgRouting functions:

\end{enumerate}

createdb -U postgres osm
createlang -U postgres plpgsql osm

psql -U postgres -f /usr/share/postgresql-8.3-postgis/lwpostgis.sql osm
psql -U postgres -f /usr/share/postgresql-8.3-postgis/spatial\_ref\_sys.sql osm

psql -U postgres -f /usr/share/postlbs/routing\_core.sql osm
psql -U postgres -f /usr/share/postlbs/routing\_core\_wrappers.sql osm
psql -U postgres -f /usr/share/postlbs/routing\_topology.sql osm
2. You can define the features and attributes to be imported from the OpenStreetMap XML file in the configuration file (default: mapconfig.xml)
\begin{enumerate}
\item {} 
Open a terminal window and run osm2pgrouting with the following paramters

\end{enumerate}
\begin{description}
\item[./osm2pgrouting -file /home/foss4g/capetown\_20080829.osm ] \leavevmode
-conf mapconfig.xml -dbname osm -user postgres -clean

\end{description}

Other available parameters are:
\begin{itemize}
\item {} \begin{description}
\item[required:] \leavevmode\begin{optionlist}{3cm}
\item [-file]  
\textless{}file\textgreater{}          -- name of your osm xml file
\end{optionlist}

-dbname \textless{}dbname\textgreater{}        -- name of your database
-user   \textless{}user\textgreater{}          -- name of the user, which have write access to the database
-conf   \textless{}file\textgreater{}          -- name of your configuration xml file

\end{description}

\item {} \begin{description}
\item[optional:] \leavevmode\begin{optionlist}{3cm}
\item [-host]  
\textless{}host\textgreater{}          -- host of your postgresql database (default: 127.0.0.1)
\item [-port]  
\textless{}port\textgreater{}          -- port of your database (default: 5432)
\end{optionlist}

-passwd \textless{}passwd\textgreater{}        --  password for database access
-clean                          -- drop peviously created tables

\end{description}

\end{itemize}
\begin{enumerate}
\item {} 
Connect to your database and see the tables that have been created

\end{enumerate}

psql -U postgres osm
d
\begin{quote}
\begin{quote}

List of relations
\end{quote}

Schema \textbar{}        Name         \textbar{}   Type   \textbar{}  Owner
\end{quote}
\begin{description}
\item[--------+---------------------+----------+----------] \leavevmode
public \textbar{} classes             \textbar{} table    \textbar{} postgres
public \textbar{} geometry\_columns    \textbar{} table    \textbar{} postgres
public \textbar{} nodes               \textbar{} table    \textbar{} postgres
public \textbar{} spatial\_ref\_sys     \textbar{} table    \textbar{} postgres
public \textbar{} types               \textbar{} table    \textbar{} postgres
public \textbar{} vertices\_tmp        \textbar{} table    \textbar{} postgres
public \textbar{} vertices\_tmp\_id\_seq \textbar{} sequence \textbar{} postgres
public \textbar{} ways                \textbar{} table    \textbar{} postgres

\end{description}

(8 rows)
Note: If tables are missing you might have forgotten to add PostGIS or pgRouting functions to your database.

Let's do some more advanced routing with those extra information about road types and road classes.

\resetcurrentobjects
\hypertarget{--doc-chapters/topology}{}

\chapter{Load your network data and create a network topology}

Some network data already comes with a network topology that can be used with pgRouting immediately. But usually the data is in a different format than we need for pgRouting. Often network data is stored in the Shape file format (.shp) and we can use PostGIS' shape2postgresql converter to import the data into the database. OpenStreetMap stores its data as XML and it has its own importing tools for PostgreSQL database.

Later we will use the osm2pgrouting converter. But it does much more than the basic steps for simple routing, so we will start this workshop with the minimum required attributes.


\section{Load the network data}

After creating the workshop database and adding the PostGIS and pgRouting functions to this database (see previous chapter), we load the sample data to our database:

\begin{Verbatim}[commandchars=@\[\]]
psql -U postgres routing
\i /home/foss4g/ways@_without@_topology.sql
\end{Verbatim}

\begin{notice}{note}{Note:}
The SQL dump file was made from a database which already had PostGIS functions loaded, so it will report errors during import that these functions already exist. You can ignore these errors.
\end{notice}

Let's see witch tables have been created:

\begin{Verbatim}[commandchars=@\[\]]
\d

                  List of relations
 Schema @textbar[]       Name       @textbar[] Type  @textbar[]  Owner
--------+------------------+-------+----------
 public @textbar[] geometry@_columns @textbar[] table @textbar[] postgres
 public @textbar[] spatial@_ref@_sys  @textbar[] table @textbar[] postgres
 public @textbar[] ways             @textbar[] table @textbar[] postgres
(3 rows)
\end{Verbatim}

\begin{Verbatim}[commandchars=@\[\]]
\d ways

               Table "public.ways"
  Column  @textbar[]       Type       @textbar[] Modifiers
----------+------------------+-----------
 gid      @textbar[] integer          @textbar[] not null
 length   @textbar[] double precision @textbar[]
 name     @textbar[] character(200)   @textbar[]
 the@_geom @textbar[] geometry         @textbar[]
Indexes:
        "ways@_pkey" PRIMARY KEY, btree (gid)
Check constraints:
        "enforce@_dims@_the@_geom" CHECK (ndims(the@_geom) = 2)
        "enforce@_geotype@_the@_geom" CHECK (geometrytype(the@_geom) =
                                        'MULTILINESTRING'::text OR the@_geom IS NULL)
        "enforce@_srid@_the@_geom" CHECK (srid(the@_geom) = 4326)
Create network topology
\end{Verbatim}

Having your data imported into a PostgreSQL database usually requires one more step for pgRouting. You have to make sure that your data provides a correct network topology, which consists of links with source and target ID each.

If your network data doesn't have such network topology information already you need to run the ``assign\_vertex\_id'' function. This function assigns a source and a target ID to each link and it can ``snap'' nearby vertices within a certain tolerance.:

\begin{Verbatim}[commandchars=\\\{\}]
\PYG{n}{assign\PYGZus{}vertex\PYGZus{}id}\PYG{p}{(}\PYG{l+s+s1}{'\textless{}table\textgreater{}'}\PYG{p}{,} \PYG{n+nb}{float} \PYG{n}{tolerance}\PYG{p}{,} \PYG{l+s+s1}{'\textless{}geometry column'}\PYG{p}{,} \PYG{l+s+s1}{'\textless{}gid\textgreater{}'}\PYG{p}{)}
\end{Verbatim}

First we have to add source and target column, then we run the assign\_vertex\_id function ... and wait.:

\begin{Verbatim}[commandchars=\\\{\}]
\PYG{k}{ALTER} \PYG{k}{TABLE} \PYG{n}{ways} \PYG{k}{ADD} \PYG{k}{COLUMN} \PYG{k}{source} \PYG{n+nb}{integer}\PYG{p}{;}
\PYG{k}{ALTER} \PYG{k}{TABLE} \PYG{n}{ways} \PYG{k}{ADD} \PYG{k}{COLUMN} \PYG{n}{target} \PYG{n+nb}{integer}\PYG{p}{;}
\PYG{k}{SELECT} \PYG{n}{assign\PYGZus{}vertex\PYGZus{}id}\PYG{p}{(}\PYG{l+s+s1}{'ways'}\PYG{p}{,} \PYG{l+m+mi}{0}\PYG{p}{.}\PYG{l+m+mi}{00001}\PYG{p}{,} \PYG{l+s+s1}{'the\PYGZus{}geom'}\PYG{p}{,} \PYG{l+s+s1}{'gid'}\PYG{p}{)}\PYG{p}{;}
\end{Verbatim}

\begin{notice}{warning}{Warning:}
The dimension of the tolerance parameter depends on your data projection. Usually it's either ``degrees'' or ``meters''. Because OSM data has a very good quality for Cape town we can choose a very small ``snapping'' tolerance: 0.00001 degrees
\end{notice}


\section{Add indices}

Fortunately we didn't need to wait too long because the data is small. But your network data might be very large, so it's a good idea to add an index on source, target and geometry column.:

\begin{Verbatim}[commandchars=\\\{\}]
\PYG{k}{CREATE} \PYG{k}{INDEX} \PYG{n}{source\PYGZus{}idx} \PYG{k}{ON} \PYG{n}{ways}\PYG{p}{(}\PYG{k}{source}\PYG{p}{)}\PYG{p}{;}
\PYG{k}{CREATE} \PYG{k}{INDEX} \PYG{n}{target\PYGZus{}idx} \PYG{k}{ON} \PYG{n}{ways}\PYG{p}{(}\PYG{n}{target}\PYG{p}{)}\PYG{p}{;}
\PYG{k}{CREATE} \PYG{k}{INDEX} \PYG{n}{geom\PYGZus{}idx} \PYG{k}{ON} \PYG{n}{ways} \PYG{k}{USING} \PYG{n}{GIST}\PYG{p}{(}\PYG{n}{the\PYGZus{}geom} \PYG{n}{GIST\PYGZus{}GEOMETRY\PYGZus{}OPS}\PYG{p}{)}\PYG{p}{;}
\end{Verbatim}

After these steps our routing database look like this:

\begin{Verbatim}[commandchars=@\[\]]
\d

                     List of relations
 Schema @textbar[]        Name         @textbar[]   Type   @textbar[]  Owner
--------+---------------------+----------+----------
 public @textbar[] geometry@_columns    @textbar[] table    @textbar[] postgres
 public @textbar[] spatial@_ref@_sys     @textbar[] table    @textbar[] postgres
 public @textbar[] vertices@_tmp        @textbar[] table    @textbar[] postgres
 public @textbar[] vertices@_tmp@_id@_seq @textbar[] sequence @textbar[] postgres
 public @textbar[] ways                @textbar[] table    @textbar[] postgres
(5 rows)
\end{Verbatim}

\begin{Verbatim}[commandchars=@\[\]]
\d ways

               Table "public.ways"
  Column  @textbar[]       Type       @textbar[] Modifiers
----------+------------------+-----------
 gid      @textbar[] integer          @textbar[] not null
 length   @textbar[] double precision @textbar[]
 name     @textbar[] character(200)   @textbar[]
 the@_geom @textbar[] geometry         @textbar[]
 source   @textbar[] integer          @textbar[]
 target   @textbar[] integer          @textbar[]
Indexes:
        "ways@_pkey" PRIMARY KEY, btree (gid)
Check constraints:
        "enforce@_dims@_the@_geom" CHECK (ndims(the@_geom) = 2)
        "enforce@_geotype@_the@_geom" CHECK (geometrytype(the@_geom) =
                                        'MULTILINESTRING'::text OR the@_geom IS NULL)
        "enforce@_srid@_the@_geom" CHECK (srid(the@_geom) = 4326)
\end{Verbatim}

Now we are ready for routing with Dijkstra algorithm!

\resetcurrentobjects
\hypertarget{--doc-chapters/shortest_path}{}

\chapter{Shortest Path Search}
\hypertarget{todo-37}{}
\begin{notice}{note}{Todo}

Add chapter introduction for Shortest Path Search
\end{notice}


\section{Dijkstra algorithm}

Dijkstra algorithm was the first algorithm implemented in pgRouting. It doesn't require more attributes than source and target ID, and it can distinguish between directed and undirected graphs. You can specify if your network has ``reverse cost'' or not.

\begin{Verbatim}[commandchars=\\\{\}]
\PYG{n}{shortest\PYGZus{}path}\PYG{p}{(} \PYG{k}{sql} \PYG{n+nb}{text}\PYG{p}{,}
                   \PYG{n}{source\PYGZus{}id} \PYG{n+nb}{integer}\PYG{p}{,}
                   \PYG{n}{target\PYGZus{}id} \PYG{n+nb}{integer}\PYG{p}{,}
                   \PYG{n}{directed} \PYG{n+nb}{boolean}\PYG{p}{,}
                   \PYG{n}{has\PYGZus{}reverse\PYGZus{}cost} \PYG{n+nb}{boolean} \PYG{p}{)}
\end{Verbatim}

\begin{notice}{note}{Note:}\begin{itemize}
\item {} 
Source and target IDs are vertex IDs.

\item {} 
Undirected graphs (``directed false'') ignores ``has\_reverse\_cost'' setting

\item {} 
Shortest Path Dijkstra core function

\end{itemize}
\end{notice}


\subsection{Core}

Each algorithm has its core function (implementation), which is the base for its wrapper functions.

\begin{Verbatim}[commandchars=\\\{\}]
\PYG{k}{SELECT} \PYG{o}{*} \PYG{k}{FROM} \PYG{n}{shortest\PYGZus{}path}\PYG{p}{(}\PYG{l+s+s1}{'}
\PYG{l+s+s1}{                SELECT gid as id,}
\PYG{l+s+s1}{                         source::integer,}
\PYG{l+s+s1}{                         target::integer,}
\PYG{l+s+s1}{                         length::double precision as cost}
\PYG{l+s+s1}{                        FROM ways'}\PYG{p}{,}
                \PYG{l+m+mi}{10}\PYG{p}{,} \PYG{l+m+mi}{20}\PYG{p}{,} \PYG{k}{false}\PYG{p}{,} \PYG{k}{false}\PYG{p}{)}\PYG{p}{;}
\end{Verbatim}

\begin{Verbatim}[commandchars=\\\{\}]
 \PYG{n}{vertex\PYGZus{}id} \PYG{o}{\textbar{}} \PYG{n}{edge\PYGZus{}id} \PYG{o}{\textbar{}}        \PYG{n}{cost}
\PYG{c+c1}{-----------+---------+---------------------}
        \PYG{l+m+mi}{10} \PYG{o}{\textbar{}}     \PYG{l+m+mi}{293} \PYG{o}{\textbar{}}  \PYG{l+m+mi}{0}\PYG{p}{.}\PYG{l+m+mi}{0059596293824534}
         \PYG{l+m+mi}{9} \PYG{o}{\textbar{}}    \PYG{l+m+mi}{4632} \PYG{o}{\textbar{}}  \PYG{l+m+mi}{0}\PYG{p}{.}\PYG{l+m+mi}{0846731039249787}
      \PYG{l+m+mi}{3974} \PYG{o}{\textbar{}}    \PYG{l+m+mi}{4633} \PYG{o}{\textbar{}}  \PYG{l+m+mi}{0}\PYG{p}{.}\PYG{l+m+mi}{0765635090514303}
      \PYG{l+m+mi}{2107} \PYG{o}{\textbar{}}    \PYG{l+m+mi}{4634} \PYG{o}{\textbar{}}  \PYG{l+m+mi}{0}\PYG{p}{.}\PYG{l+m+mi}{0763951531894937}
       \PYG{p}{.}\PYG{p}{.}\PYG{p}{.} \PYG{o}{\textbar{}}     \PYG{p}{.}\PYG{p}{.}\PYG{p}{.} \PYG{o}{\textbar{}}  \PYG{p}{.}\PYG{p}{.}\PYG{p}{.}
        \PYG{l+m+mi}{20} \PYG{o}{\textbar{}}      \PYG{o}{-}\PYG{l+m+mi}{1} \PYG{o}{\textbar{}}                   \PYG{l+m+mi}{0}
\PYG{p}{(}\PYG{l+m+mi}{63} \PYG{k}{rows}\PYG{p}{)}
\end{Verbatim}


\subsection{Wrapper}
\paragraph{Wrapper WITHOUT bounding box}

Wrapper functions extend the core functions with transformations, bounding box limitations, etc.. Wrappers can change the format and ordering of the result. They often set default function parameters and make the usage of pgRouting more simple.

\begin{Verbatim}[commandchars=\\\{\}]
\PYG{k}{SELECT} \PYG{n}{gid}\PYG{p}{,} \PYG{n}{AsText}\PYG{p}{(}\PYG{n}{the\PYGZus{}geom}\PYG{p}{)} \PYG{k}{AS} \PYG{n}{the\PYGZus{}geom}
        \PYG{k}{FROM} \PYG{n}{dijkstra\PYGZus{}sp}\PYG{p}{(}\PYG{l+s+s1}{'ways'}\PYG{p}{,} \PYG{l+m+mi}{10}\PYG{p}{,} \PYG{l+m+mi}{20}\PYG{p}{)}\PYG{p}{;}
\end{Verbatim}

\begin{Verbatim}[commandchars=\\\{\}]
  \PYG{n}{gid}   \PYG{o}{\textbar{}}                              \PYG{n}{the\PYGZus{}geom}
\PYG{c+c1}{--------+---------------------------------------------------------------}
    \PYG{l+m+mi}{293} \PYG{o}{\textbar{}} \PYG{n}{MULTILINESTRING}\PYG{p}{(}\PYG{p}{(}\PYG{l+m+mi}{18}\PYG{p}{.}\PYG{l+m+mi}{4074149} \PYG{o}{-}\PYG{l+m+mi}{33}\PYG{p}{.}\PYG{l+m+mi}{9443308}\PYG{p}{,}\PYG{l+m+mi}{18}\PYG{p}{.}\PYG{l+m+mi}{4074019} \PYG{o}{-}\PYG{l+m+mi}{33}\PYG{p}{.}\PYG{l+m+mi}{9443833}\PYG{p}{)}\PYG{p}{)}
   \PYG{l+m+mi}{4632} \PYG{o}{\textbar{}} \PYG{n}{MULTILINESTRING}\PYG{p}{(}\PYG{p}{(}\PYG{l+m+mi}{18}\PYG{p}{.}\PYG{l+m+mi}{4074149} \PYG{o}{-}\PYG{l+m+mi}{33}\PYG{p}{.}\PYG{l+m+mi}{9443308}\PYG{p}{,}\PYG{l+m+mi}{18}\PYG{p}{.}\PYG{l+m+mi}{4077388} \PYG{o}{-}\PYG{l+m+mi}{33}\PYG{p}{.}\PYG{l+m+mi}{9436183}\PYG{p}{)}\PYG{p}{)}
   \PYG{l+m+mi}{4633} \PYG{o}{\textbar{}} \PYG{n}{MULTILINESTRING}\PYG{p}{(}\PYG{p}{(}\PYG{l+m+mi}{18}\PYG{p}{.}\PYG{l+m+mi}{4077388} \PYG{o}{-}\PYG{l+m+mi}{33}\PYG{p}{.}\PYG{l+m+mi}{9436183}\PYG{p}{,}\PYG{l+m+mi}{18}\PYG{p}{.}\PYG{l+m+mi}{4080293} \PYG{o}{-}\PYG{l+m+mi}{33}\PYG{p}{.}\PYG{l+m+mi}{9429733}\PYG{p}{)}\PYG{p}{)}
    \PYG{p}{.}\PYG{p}{.}\PYG{p}{.} \PYG{o}{\textbar{}} \PYG{p}{.}\PYG{p}{.}\PYG{p}{.}
    \PYG{l+m+mi}{762} \PYG{o}{\textbar{}} \PYG{n}{MULTILINESTRING}\PYG{p}{(}\PYG{p}{(}\PYG{l+m+mi}{18}\PYG{p}{.}\PYG{l+m+mi}{4241422} \PYG{o}{-}\PYG{l+m+mi}{33}\PYG{p}{.}\PYG{l+m+mi}{9179275}\PYG{p}{,}\PYG{l+m+mi}{18}\PYG{p}{.}\PYG{l+m+mi}{4237423} \PYG{o}{-}\PYG{l+m+mi}{33}\PYG{p}{.}\PYG{l+m+mi}{9182966}\PYG{p}{)}\PYG{p}{)}
    \PYG{l+m+mi}{761} \PYG{o}{\textbar{}} \PYG{n}{MULTILINESTRING}\PYG{p}{(}\PYG{p}{(}\PYG{l+m+mi}{18}\PYG{p}{.}\PYG{l+m+mi}{4243523} \PYG{o}{-}\PYG{l+m+mi}{33}\PYG{p}{.}\PYG{l+m+mi}{9177154}\PYG{p}{,}\PYG{l+m+mi}{18}\PYG{p}{.}\PYG{l+m+mi}{4241422} \PYG{o}{-}\PYG{l+m+mi}{33}\PYG{p}{.}\PYG{l+m+mi}{9179275}\PYG{p}{)}\PYG{p}{)}
\PYG{p}{(}\PYG{l+m+mi}{62} \PYG{k}{rows}\PYG{p}{)}
\end{Verbatim}
\paragraph{Wrapper WITH bounding box}

You can limit your search area by adding a bounding box. This will improve performance especially for large networks.

\begin{Verbatim}[commandchars=\\\{\}]
\PYG{k}{SELECT} \PYG{n}{gid}\PYG{p}{,} \PYG{n}{AsText}\PYG{p}{(}\PYG{n}{the\PYGZus{}geom}\PYG{p}{)} \PYG{k}{AS} \PYG{n}{the\PYGZus{}geom}
        \PYG{k}{FROM} \PYG{n}{dijkstra\PYGZus{}sp\PYGZus{}delta}\PYG{p}{(}\PYG{l+s+s1}{'ways'}\PYG{p}{,} \PYG{l+m+mi}{10}\PYG{p}{,} \PYG{l+m+mi}{20}\PYG{p}{,} \PYG{l+m+mi}{0}\PYG{p}{.}\PYG{l+m+mi}{1}\PYG{p}{)}\PYG{p}{;}
\end{Verbatim}

\begin{Verbatim}[commandchars=\\\{\}]
   \PYG{n}{gid}  \PYG{o}{\textbar{}} \PYG{n}{the\PYGZus{}geom}
\PYG{c+c1}{--------+---------------------------------------------------------------}
   \PYG{l+m+mi}{293}  \PYG{o}{\textbar{}} \PYG{n}{MULTILINESTRING}\PYG{p}{(}\PYG{p}{(}\PYG{l+m+mi}{18}\PYG{p}{.}\PYG{l+m+mi}{4074149} \PYG{o}{-}\PYG{l+m+mi}{33}\PYG{p}{.}\PYG{l+m+mi}{9443308}\PYG{p}{,}\PYG{l+m+mi}{18}\PYG{p}{.}\PYG{l+m+mi}{4074019} \PYG{o}{-}\PYG{l+m+mi}{33}\PYG{p}{.}\PYG{l+m+mi}{9443833}\PYG{p}{)}\PYG{p}{)}
   \PYG{l+m+mi}{4632} \PYG{o}{\textbar{}} \PYG{n}{MULTILINESTRING}\PYG{p}{(}\PYG{p}{(}\PYG{l+m+mi}{18}\PYG{p}{.}\PYG{l+m+mi}{4074149} \PYG{o}{-}\PYG{l+m+mi}{33}\PYG{p}{.}\PYG{l+m+mi}{9443308}\PYG{p}{,}\PYG{l+m+mi}{18}\PYG{p}{.}\PYG{l+m+mi}{4077388} \PYG{o}{-}\PYG{l+m+mi}{33}\PYG{p}{.}\PYG{l+m+mi}{9436183}\PYG{p}{)}\PYG{p}{)}
   \PYG{l+m+mi}{4633} \PYG{o}{\textbar{}} \PYG{n}{MULTILINESTRING}\PYG{p}{(}\PYG{p}{(}\PYG{l+m+mi}{18}\PYG{p}{.}\PYG{l+m+mi}{4077388} \PYG{o}{-}\PYG{l+m+mi}{33}\PYG{p}{.}\PYG{l+m+mi}{9436183}\PYG{p}{,}\PYG{l+m+mi}{18}\PYG{p}{.}\PYG{l+m+mi}{4080293} \PYG{o}{-}\PYG{l+m+mi}{33}\PYG{p}{.}\PYG{l+m+mi}{9429733}\PYG{p}{)}\PYG{p}{)}
   \PYG{p}{.}\PYG{p}{.}\PYG{p}{.}  \PYG{o}{\textbar{}} \PYG{p}{.}\PYG{p}{.}\PYG{p}{.}
   \PYG{l+m+mi}{762}  \PYG{o}{\textbar{}} \PYG{n}{MULTILINESTRING}\PYG{p}{(}\PYG{p}{(}\PYG{l+m+mi}{18}\PYG{p}{.}\PYG{l+m+mi}{4241422} \PYG{o}{-}\PYG{l+m+mi}{33}\PYG{p}{.}\PYG{l+m+mi}{9179275}\PYG{p}{,}\PYG{l+m+mi}{18}\PYG{p}{.}\PYG{l+m+mi}{4237423} \PYG{o}{-}\PYG{l+m+mi}{33}\PYG{p}{.}\PYG{l+m+mi}{9182966}\PYG{p}{)}\PYG{p}{)}
   \PYG{l+m+mi}{761} \PYG{o}{\textbar{}} \PYG{n}{MULTILINESTRING}\PYG{p}{(}\PYG{p}{(}\PYG{l+m+mi}{18}\PYG{p}{.}\PYG{l+m+mi}{4243523} \PYG{o}{-}\PYG{l+m+mi}{33}\PYG{p}{.}\PYG{l+m+mi}{9177154}\PYG{p}{,}\PYG{l+m+mi}{18}\PYG{p}{.}\PYG{l+m+mi}{4241422} \PYG{o}{-}\PYG{l+m+mi}{33}\PYG{p}{.}\PYG{l+m+mi}{9179275}\PYG{p}{)}\PYG{p}{)}
\PYG{p}{(}\PYG{l+m+mi}{62} \PYG{k}{rows}\PYG{p}{)}
\end{Verbatim}

\begin{notice}{warning}{Warning:}
The projection of OSM data is ``degree'', so we set a bounding box containing start and end vertex plus a 0.1 degree buffer for example.
\end{notice}


\section{A-Star algorithm}

A-Star algorithm is another well-known routing algorithm. It adds geographical information to source and target of each network link. This enables the shortest path search to prefer links which are closer to the target of the search.


\subsection{Prerequisites}

For A-Star you need to prepare your network table and add latitute/longitude columns (x1, y1 and x2, y2) and calculate their values.

\begin{Verbatim}[commandchars=\\\{\}]
\PYG{k}{ALTER} \PYG{k}{TABLE} \PYG{n}{ways} \PYG{k}{ADD} \PYG{k}{COLUMN} \PYG{n}{x1} \PYG{n}{double} \PYG{k}{precision}\PYG{p}{;}
\PYG{k}{ALTER} \PYG{k}{TABLE} \PYG{n}{ways} \PYG{k}{ADD} \PYG{k}{COLUMN} \PYG{n}{y1} \PYG{n}{double} \PYG{k}{precision}\PYG{p}{;}
\PYG{k}{ALTER} \PYG{k}{TABLE} \PYG{n}{ways} \PYG{k}{ADD} \PYG{k}{COLUMN} \PYG{n}{x2} \PYG{n}{double} \PYG{k}{precision}\PYG{p}{;}
\PYG{k}{ALTER} \PYG{k}{TABLE} \PYG{n}{ways} \PYG{k}{ADD} \PYG{k}{COLUMN} \PYG{n}{y2} \PYG{n}{double} \PYG{k}{precision}\PYG{p}{;}

\PYG{k}{UPDATE} \PYG{n}{ways} \PYG{k}{SET} \PYG{n}{x1} \PYG{o}{=} \PYG{n}{x}\PYG{p}{(}\PYG{n}{startpoint}\PYG{p}{(}\PYG{n}{the\PYGZus{}geom}\PYG{p}{)}\PYG{p}{)}\PYG{p}{;}
\PYG{k}{UPDATE} \PYG{n}{ways} \PYG{k}{SET} \PYG{n}{y1} \PYG{o}{=} \PYG{n}{y}\PYG{p}{(}\PYG{n}{startpoint}\PYG{p}{(}\PYG{n}{the\PYGZus{}geom}\PYG{p}{)}\PYG{p}{)}\PYG{p}{;}

\PYG{k}{UPDATE} \PYG{n}{ways} \PYG{k}{SET} \PYG{n}{x2} \PYG{o}{=} \PYG{n}{x}\PYG{p}{(}\PYG{n}{endpoint}\PYG{p}{(}\PYG{n}{the\PYGZus{}geom}\PYG{p}{)}\PYG{p}{)}\PYG{p}{;}
\PYG{k}{UPDATE} \PYG{n}{ways} \PYG{k}{SET} \PYG{n}{y2} \PYG{o}{=} \PYG{n}{y}\PYG{p}{(}\PYG{n}{endpoint}\PYG{p}{(}\PYG{n}{the\PYGZus{}geom}\PYG{p}{)}\PYG{p}{)}\PYG{p}{;}

\PYG{k}{UPDATE} \PYG{n}{ways} \PYG{k}{SET} \PYG{n}{x1} \PYG{o}{=} \PYG{n}{x}\PYG{p}{(}\PYG{n}{PointN}\PYG{p}{(}\PYG{n}{the\PYGZus{}geom}\PYG{p}{,} \PYG{l+m+mi}{1}\PYG{p}{)}\PYG{p}{)}\PYG{p}{;}
\PYG{k}{UPDATE} \PYG{n}{ways} \PYG{k}{SET} \PYG{n}{y1} \PYG{o}{=} \PYG{n}{y}\PYG{p}{(}\PYG{n}{PointN}\PYG{p}{(}\PYG{n}{the\PYGZus{}geom}\PYG{p}{,} \PYG{l+m+mi}{1}\PYG{p}{)}\PYG{p}{)}\PYG{p}{;}

\PYG{k}{UPDATE} \PYG{n}{ways} \PYG{k}{SET} \PYG{n}{x2} \PYG{o}{=} \PYG{n}{x}\PYG{p}{(}\PYG{n}{PointN}\PYG{p}{(}\PYG{n}{the\PYGZus{}geom}\PYG{p}{,} \PYG{n}{NumPoints}\PYG{p}{(}\PYG{n}{the\PYGZus{}geom}\PYG{p}{)}\PYG{p}{)}\PYG{p}{)}\PYG{p}{;}
\PYG{k}{UPDATE} \PYG{n}{ways} \PYG{k}{SET} \PYG{n}{y2} \PYG{o}{=} \PYG{n}{y}\PYG{p}{(}\PYG{n}{PointN}\PYG{p}{(}\PYG{n}{the\PYGZus{}geom}\PYG{p}{,} \PYG{n}{NumPoints}\PYG{p}{(}\PYG{n}{the\PYGZus{}geom}\PYG{p}{)}\PYG{p}{)}\PYG{p}{)}\PYG{p}{;}
\end{Verbatim}

\begin{notice}{note}{Note:}
``endpoint()'' function fails for some versions of PostgreSQL (ie. 8.2.5, 8.1.9). A workaround for that problem is using the ``PointN()'' function instead:
\end{notice}


\subsection{Core}

Shortest Path A-Star function is very similar to the Dijkstra function, though it prefers links that are close to the target of the search. The heuristics of this search are predefined, so you need to recompile pgRouting if you want to make changes to the heuristic function itself.

\begin{Verbatim}[commandchars=\\\{\}]
\PYG{n}{shortest\PYGZus{}path\PYGZus{}astar}\PYG{p}{(} \PYG{k}{sql} \PYG{n+nb}{text}\PYG{p}{,}
                   \PYG{n}{source\PYGZus{}id} \PYG{n+nb}{integer}\PYG{p}{,}
                   \PYG{n}{target\PYGZus{}id} \PYG{n+nb}{integer}\PYG{p}{,}
                   \PYG{n}{directed} \PYG{n+nb}{boolean}\PYG{p}{,}
                   \PYG{n}{has\PYGZus{}reverse\PYGZus{}cost} \PYG{n+nb}{boolean} \PYG{p}{)}
\end{Verbatim}

\begin{notice}{note}{Note:}\begin{itemize}
\item {} 
Source and target IDs are vertex IDs.

\item {} 
Undirected graphs (``directed false'') ignores ``has\_reverse\_cost'' setting

\item {} 
Example of A-Star core function

\end{itemize}
\end{notice}

\begin{Verbatim}[commandchars=\\\{\}]
\PYG{k}{SELECT} \PYG{o}{*} \PYG{k}{FROM} \PYG{n}{shortest\PYGZus{}path\PYGZus{}astar}\PYG{p}{(}\PYG{l+s+s1}{'}
\PYG{l+s+s1}{                SELECT gid as id,}
\PYG{l+s+s1}{                         source::integer,}
\PYG{l+s+s1}{                         target::integer,}
\PYG{l+s+s1}{                         length::double precision as cost,}
\PYG{l+s+s1}{                         x1, y1, x2, y2}
\PYG{l+s+s1}{                        FROM ways'}\PYG{p}{,}
                \PYG{l+m+mi}{10}\PYG{p}{,} \PYG{l+m+mi}{20}\PYG{p}{,} \PYG{k}{false}\PYG{p}{,} \PYG{k}{false}\PYG{p}{)}\PYG{p}{;}
\end{Verbatim}

\begin{Verbatim}[commandchars=\\\{\}]
\PYG{n}{vertex\PYGZus{}id} \PYG{o}{\textbar{}} \PYG{n}{edge\PYGZus{}id} \PYG{o}{\textbar{}}        \PYG{n}{cost}
\PYG{c+c1}{-----------+---------+---------------------}
       \PYG{l+m+mi}{10} \PYG{o}{\textbar{}}     \PYG{l+m+mi}{293} \PYG{o}{\textbar{}}  \PYG{l+m+mi}{0}\PYG{p}{.}\PYG{l+m+mi}{0059596293824534}
        \PYG{l+m+mi}{9} \PYG{o}{\textbar{}}    \PYG{l+m+mi}{4632} \PYG{o}{\textbar{}}  \PYG{l+m+mi}{0}\PYG{p}{.}\PYG{l+m+mi}{0846731039249787}
     \PYG{l+m+mi}{3974} \PYG{o}{\textbar{}}    \PYG{l+m+mi}{4633} \PYG{o}{\textbar{}}  \PYG{l+m+mi}{0}\PYG{p}{.}\PYG{l+m+mi}{0765635090514303}
      \PYG{p}{.}\PYG{p}{.}\PYG{p}{.} \PYG{o}{\textbar{}}     \PYG{p}{.}\PYG{p}{.}\PYG{p}{.} \PYG{o}{\textbar{}}  \PYG{p}{.}\PYG{p}{.}\PYG{p}{.}
       \PYG{l+m+mi}{20} \PYG{o}{\textbar{}}      \PYG{o}{-}\PYG{l+m+mi}{1} \PYG{o}{\textbar{}}                   \PYG{l+m+mi}{0}
\PYG{p}{(}\PYG{l+m+mi}{63} \PYG{k}{rows}\PYG{p}{)}
\end{Verbatim}


\subsection{Wrapper}
\paragraph{Wrapper function WITH bounding box}

Wrapper functions extend the core functions with transformations, bounding box limitations, etc..

\begin{Verbatim}[commandchars=\\\{\}]
\PYG{k}{SELECT} \PYG{n}{gid}\PYG{p}{,} \PYG{n}{AsText}\PYG{p}{(}\PYG{n}{the\PYGZus{}geom}\PYG{p}{)} \PYG{k}{AS} \PYG{n}{the\PYGZus{}geom}
        \PYG{k}{FROM} \PYG{n}{astar\PYGZus{}sp\PYGZus{}delta}\PYG{p}{(}\PYG{l+s+s1}{'ways'}\PYG{p}{,} \PYG{l+m+mi}{10}\PYG{p}{,} \PYG{l+m+mi}{20}\PYG{p}{,} \PYG{l+m+mi}{0}\PYG{p}{.}\PYG{l+m+mi}{1}\PYG{p}{)}\PYG{p}{;}
\end{Verbatim}

\begin{Verbatim}[commandchars=\\\{\}]
  \PYG{n}{gid}   \PYG{o}{\textbar{}}                              \PYG{n}{the\PYGZus{}geom}
\PYG{c+c1}{--------+---------------------------------------------------------------}
    \PYG{l+m+mi}{293} \PYG{o}{\textbar{}} \PYG{n}{MULTILINESTRING}\PYG{p}{(}\PYG{p}{(}\PYG{l+m+mi}{18}\PYG{p}{.}\PYG{l+m+mi}{4074149} \PYG{o}{-}\PYG{l+m+mi}{33}\PYG{p}{.}\PYG{l+m+mi}{9443308}\PYG{p}{,}\PYG{l+m+mi}{18}\PYG{p}{.}\PYG{l+m+mi}{4074019} \PYG{o}{-}\PYG{l+m+mi}{33}\PYG{p}{.}\PYG{l+m+mi}{9443833}\PYG{p}{)}\PYG{p}{)}
   \PYG{l+m+mi}{4632} \PYG{o}{\textbar{}} \PYG{n}{MULTILINESTRING}\PYG{p}{(}\PYG{p}{(}\PYG{l+m+mi}{18}\PYG{p}{.}\PYG{l+m+mi}{4074149} \PYG{o}{-}\PYG{l+m+mi}{33}\PYG{p}{.}\PYG{l+m+mi}{9443308}\PYG{p}{,}\PYG{l+m+mi}{18}\PYG{p}{.}\PYG{l+m+mi}{4077388} \PYG{o}{-}\PYG{l+m+mi}{33}\PYG{p}{.}\PYG{l+m+mi}{9436183}\PYG{p}{)}\PYG{p}{)}
   \PYG{l+m+mi}{4633} \PYG{o}{\textbar{}} \PYG{n}{MULTILINESTRING}\PYG{p}{(}\PYG{p}{(}\PYG{l+m+mi}{18}\PYG{p}{.}\PYG{l+m+mi}{4077388} \PYG{o}{-}\PYG{l+m+mi}{33}\PYG{p}{.}\PYG{l+m+mi}{9436183}\PYG{p}{,}\PYG{l+m+mi}{18}\PYG{p}{.}\PYG{l+m+mi}{4080293} \PYG{o}{-}\PYG{l+m+mi}{33}\PYG{p}{.}\PYG{l+m+mi}{9429733}\PYG{p}{)}\PYG{p}{)}
    \PYG{p}{.}\PYG{p}{.}\PYG{p}{.} \PYG{o}{\textbar{}} \PYG{p}{.}\PYG{p}{.}\PYG{p}{.}
    \PYG{l+m+mi}{762} \PYG{o}{\textbar{}} \PYG{n}{MULTILINESTRING}\PYG{p}{(}\PYG{p}{(}\PYG{l+m+mi}{18}\PYG{p}{.}\PYG{l+m+mi}{4241422} \PYG{o}{-}\PYG{l+m+mi}{33}\PYG{p}{.}\PYG{l+m+mi}{9179275}\PYG{p}{,}\PYG{l+m+mi}{18}\PYG{p}{.}\PYG{l+m+mi}{4237423} \PYG{o}{-}\PYG{l+m+mi}{33}\PYG{p}{.}\PYG{l+m+mi}{9182966}\PYG{p}{)}\PYG{p}{)}
    \PYG{l+m+mi}{761} \PYG{o}{\textbar{}} \PYG{n}{MULTILINESTRING}\PYG{p}{(}\PYG{p}{(}\PYG{l+m+mi}{18}\PYG{p}{.}\PYG{l+m+mi}{4243523} \PYG{o}{-}\PYG{l+m+mi}{33}\PYG{p}{.}\PYG{l+m+mi}{9177154}\PYG{p}{,}\PYG{l+m+mi}{18}\PYG{p}{.}\PYG{l+m+mi}{4241422} \PYG{o}{-}\PYG{l+m+mi}{33}\PYG{p}{.}\PYG{l+m+mi}{9179275}\PYG{p}{)}\PYG{p}{)}
\PYG{p}{(}\PYG{l+m+mi}{62} \PYG{k}{rows}\PYG{p}{)}
\end{Verbatim}

\begin{notice}{note}{Note:}
There is currently no wrapper function for A-Star without bounding box, since bounding boxes are very useful to increase performance. If you don't need a bounding box Dijkstra will be enough anyway.
\end{notice}

\begin{notice}{warning}{Warning:}
The projection of OSM data is ``degree'', so we set a bounding box containing start and end vertex plus a 0.1 degree buffer for example.
\end{notice}


\section{Shooting-Star algorithm}

Shooting-Star algorithm is the latest of pgRouting shortest path algorithms. Its speciality is that it routes from link to link, not from vertex to vertex as Dijkstra and A-Star algorithms do. This makes it possible to define relations between links for example, and it solves some other vertex-based algorithm issues like ``parallel links'', which have same source and target but different costs.


\subsection{Prerequisites}

For Shooting-Star you need to prepare your network table and add the ``reverse\_cost'' and ``to\_cost'' column. Like A-Star this algorithm also has a heuristic function, which prefers links closer to the target of the search.

\begin{Verbatim}[commandchars=\\\{\}]
\PYG{k}{ALTER} \PYG{k}{TABLE} \PYG{n}{ways} \PYG{k}{ADD} \PYG{k}{COLUMN} \PYG{n}{reverse\PYGZus{}cost} \PYG{n}{double} \PYG{k}{precision}\PYG{p}{;}
\PYG{k}{UPDATE} \PYG{n}{ways} \PYG{k}{SET} \PYG{n}{reverse\PYGZus{}cost} \PYG{o}{=} \PYG{k}{length}\PYG{p}{;}

\PYG{k}{ALTER} \PYG{k}{TABLE} \PYG{n}{ways} \PYG{k}{ADD} \PYG{k}{COLUMN} \PYG{n}{to\PYGZus{}cost} \PYG{n}{double} \PYG{k}{precision}\PYG{p}{;}

\PYG{k}{ALTER} \PYG{k}{TABLE} \PYG{n}{ways} \PYG{k}{ADD} \PYG{k}{COLUMN} \PYG{k}{rule} \PYG{n+nb}{text}\PYG{p}{;}
\end{Verbatim}
\paragraph{Shooting-Star algorithm introduces two new attributes}
\begin{itemize}
\item {} 
\textbf{rule}: a string with a comma separated list of edge IDs, which describes a rule for turning restriction (if you came along these edges, you can pass through the current one only with the cost stated in to\_cost column)

\item {} 
\textbf{to\_cost}: a cost of a restricted passage (can be very high in a case of turn restriction or comparable with an edge cost in a case of traffic light)

\end{itemize}

\begin{Verbatim}[commandchars=\\\{\}]
\PYG{n}{shortest\PYGZus{}path\PYGZus{}shooting\PYGZus{}star}\PYG{p}{(} \PYG{k}{sql} \PYG{n+nb}{text}\PYG{p}{,}
                   \PYG{n}{source\PYGZus{}id} \PYG{n+nb}{integer}\PYG{p}{,}
                   \PYG{n}{target\PYGZus{}id} \PYG{n+nb}{integer}\PYG{p}{,}
                   \PYG{n}{directed} \PYG{n+nb}{boolean}\PYG{p}{,}
                   \PYG{n}{has\PYGZus{}reverse\PYGZus{}cost} \PYG{n+nb}{boolean} \PYG{p}{)}
\end{Verbatim}

\begin{notice}{note}{Note:}\begin{itemize}
\item {} 
Source and target IDs are link IDs.

\item {} 
Undirected graphs (``directed false'') ignores ``has\_reverse\_cost'' setting

\item {} 
Example for Shooting-Star ``rule''

\end{itemize}
\end{notice}

\begin{notice}{warning}{Warning:}
Shooting* algorithm calculates a path from edge to edge (not from vertex to vertex). Column vertex\_id contains start vertex of an edge from column edge\_id.
\end{notice}

To describe turn restrictions:

\begin{Verbatim}[commandchars=\\\{\}]
 \PYG{n}{gid} \PYG{o}{\textbar{}} \PYG{k}{source} \PYG{o}{\textbar{}} \PYG{n}{target} \PYG{o}{\textbar{}} \PYG{n}{cost} \PYG{o}{\textbar{}} \PYG{n}{x1} \PYG{o}{\textbar{}} \PYG{n}{y1} \PYG{o}{\textbar{}} \PYG{n}{x2} \PYG{o}{\textbar{}} \PYG{n}{y2} \PYG{o}{\textbar{}} \PYG{n}{to\PYGZus{}cost} \PYG{o}{\textbar{}} \PYG{k}{rule}
\PYG{c+c1}{-----+--------+--------+------+----+----+----+----+---------+------}
  \PYG{l+m+mi}{12} \PYG{o}{\textbar{}}      \PYG{l+m+mi}{3} \PYG{o}{\textbar{}}     \PYG{l+m+mi}{10} \PYG{o}{\textbar{}}    \PYG{l+m+mi}{2} \PYG{o}{\textbar{}}  \PYG{l+m+mi}{4} \PYG{o}{\textbar{}}  \PYG{l+m+mi}{3} \PYG{o}{\textbar{}}  \PYG{l+m+mi}{4} \PYG{o}{\textbar{}}  \PYG{l+m+mi}{5} \PYG{o}{\textbar{}}    \PYG{l+m+mi}{1000} \PYG{o}{\textbar{}} \PYG{l+m+mi}{14}
\end{Verbatim}

... means that the cost of going from edge 14 to edge 12 is 1000, and

\begin{Verbatim}[commandchars=\\\{\}]
 \PYG{n}{gid} \PYG{o}{\textbar{}} \PYG{k}{source} \PYG{o}{\textbar{}} \PYG{n}{target} \PYG{o}{\textbar{}} \PYG{n}{cost} \PYG{o}{\textbar{}} \PYG{n}{x1} \PYG{o}{\textbar{}} \PYG{n}{y1} \PYG{o}{\textbar{}} \PYG{n}{x2} \PYG{o}{\textbar{}} \PYG{n}{y2} \PYG{o}{\textbar{}} \PYG{n}{to\PYGZus{}cost} \PYG{o}{\textbar{}} \PYG{k}{rule}
\PYG{c+c1}{-----+--------+--------+------+----+----+----+----+---------+------}
  \PYG{l+m+mi}{12} \PYG{o}{\textbar{}}      \PYG{l+m+mi}{3} \PYG{o}{\textbar{}}     \PYG{l+m+mi}{10} \PYG{o}{\textbar{}}    \PYG{l+m+mi}{2} \PYG{o}{\textbar{}}  \PYG{l+m+mi}{4} \PYG{o}{\textbar{}}  \PYG{l+m+mi}{3} \PYG{o}{\textbar{}}  \PYG{l+m+mi}{4} \PYG{o}{\textbar{}}  \PYG{l+m+mi}{5} \PYG{o}{\textbar{}}    \PYG{l+m+mi}{1000} \PYG{o}{\textbar{}} \PYG{l+m+mi}{14}\PYG{p}{,} \PYG{l+m+mi}{4}
\end{Verbatim}

... means that the cost of going from edge 14 to edge 12 through edge 4 is 1000.

If you need multiple restrictions for a given edge then you have to add multiple records for that edge each with a separate restriction.

\begin{Verbatim}[commandchars=\\\{\}]
 \PYG{n}{gid} \PYG{o}{\textbar{}} \PYG{k}{source} \PYG{o}{\textbar{}} \PYG{n}{target} \PYG{o}{\textbar{}} \PYG{n}{cost} \PYG{o}{\textbar{}} \PYG{n}{x1} \PYG{o}{\textbar{}} \PYG{n}{y1} \PYG{o}{\textbar{}} \PYG{n}{x2} \PYG{o}{\textbar{}} \PYG{n}{y2} \PYG{o}{\textbar{}} \PYG{n}{to\PYGZus{}cost} \PYG{o}{\textbar{}} \PYG{k}{rule}
\PYG{c+c1}{-----+--------+--------+------+----+----+----+----+---------+------}
  \PYG{l+m+mi}{11} \PYG{o}{\textbar{}}      \PYG{l+m+mi}{3} \PYG{o}{\textbar{}}     \PYG{l+m+mi}{10} \PYG{o}{\textbar{}}    \PYG{l+m+mi}{2} \PYG{o}{\textbar{}}  \PYG{l+m+mi}{4} \PYG{o}{\textbar{}}  \PYG{l+m+mi}{3} \PYG{o}{\textbar{}}  \PYG{l+m+mi}{4} \PYG{o}{\textbar{}}  \PYG{l+m+mi}{5} \PYG{o}{\textbar{}}    \PYG{l+m+mi}{1000} \PYG{o}{\textbar{}} \PYG{l+m+mi}{4}
  \PYG{l+m+mi}{11} \PYG{o}{\textbar{}}      \PYG{l+m+mi}{3} \PYG{o}{\textbar{}}     \PYG{l+m+mi}{10} \PYG{o}{\textbar{}}    \PYG{l+m+mi}{2} \PYG{o}{\textbar{}}  \PYG{l+m+mi}{4} \PYG{o}{\textbar{}}  \PYG{l+m+mi}{3} \PYG{o}{\textbar{}}  \PYG{l+m+mi}{4} \PYG{o}{\textbar{}}  \PYG{l+m+mi}{5} \PYG{o}{\textbar{}}    \PYG{l+m+mi}{1000} \PYG{o}{\textbar{}} \PYG{l+m+mi}{12}
\end{Verbatim}

... means that the cost of going from either edge 4 or 12 to edge 11 is 1000. And then you always need to order your data by gid when you load it to a shortest path function..


\subsection{Core}

\begin{Verbatim}[commandchars=\\\{\}]
\PYG{k}{SELECT} \PYG{o}{*} \PYG{k}{FROM} \PYG{n}{shortest\PYGZus{}path\PYGZus{}shooting\PYGZus{}star}\PYG{p}{(}\PYG{l+s+s1}{'}
\PYG{l+s+s1}{                SELECT gid as id,}
\PYG{l+s+s1}{                         source::integer,}
\PYG{l+s+s1}{                         target::integer,}
\PYG{l+s+s1}{                         length::double precision as cost,}
\PYG{l+s+s1}{                         x1, y1, x2, y2,}
\PYG{l+s+s1}{                         rule, to\PYGZus{}cost}
\PYG{l+s+s1}{                        FROM ways'}\PYG{p}{,}
                \PYG{l+m+mi}{293}\PYG{p}{,} \PYG{l+m+mi}{761}\PYG{p}{,} \PYG{k}{false}\PYG{p}{,} \PYG{k}{false}\PYG{p}{)}\PYG{p}{;}
\end{Verbatim}

\begin{Verbatim}[commandchars=\\\{\}]
 \PYG{n}{vertex\PYGZus{}id} \PYG{o}{\textbar{}} \PYG{n}{edge\PYGZus{}id} \PYG{o}{\textbar{}}        \PYG{n}{cost}
\PYG{c+c1}{-----------+---------+---------------------}
      \PYG{l+m+mi}{4232} \PYG{o}{\textbar{}}     \PYG{l+m+mi}{293} \PYG{o}{\textbar{}}  \PYG{l+m+mi}{0}\PYG{p}{.}\PYG{l+m+mi}{0059596293824534}
      \PYG{l+m+mi}{3144} \PYG{o}{\textbar{}}     \PYG{l+m+mi}{293} \PYG{o}{\textbar{}}  \PYG{l+m+mi}{0}\PYG{p}{.}\PYG{l+m+mi}{0059596293824534}
      \PYG{l+m+mi}{4232} \PYG{o}{\textbar{}}    \PYG{l+m+mi}{4632} \PYG{o}{\textbar{}}  \PYG{l+m+mi}{0}\PYG{p}{.}\PYG{l+m+mi}{0846731039249787}
       \PYG{p}{.}\PYG{p}{.}\PYG{p}{.} \PYG{o}{\textbar{}}     \PYG{p}{.}\PYG{p}{.}\PYG{p}{.} \PYG{o}{\textbar{}}  \PYG{p}{.}\PYG{p}{.}\PYG{p}{.}
        \PYG{l+m+mi}{51} \PYG{o}{\textbar{}}     \PYG{l+m+mi}{761} \PYG{o}{\textbar{}}  \PYG{l+m+mi}{0}\PYG{p}{.}\PYG{l+m+mi}{0305298478239596}
\PYG{p}{(}\PYG{l+m+mi}{63} \PYG{k}{rows}\PYG{p}{)}
\end{Verbatim}


\subsection{Wrapper}

Wrapper functions extend the core functions with transformations, bounding box limitations, etc..

\begin{Verbatim}[commandchars=\\\{\}]
\PYG{k}{SELECT} \PYG{n}{gid}\PYG{p}{,} \PYG{n}{AsText}\PYG{p}{(}\PYG{n}{the\PYGZus{}geom}\PYG{p}{)} \PYG{k}{AS} \PYG{n}{the\PYGZus{}geom}
        \PYG{k}{FROM} \PYG{n}{shootingstar\PYGZus{}sp}\PYG{p}{(}\PYG{l+s+s1}{'ways'}\PYG{p}{,} \PYG{l+m+mi}{293}\PYG{p}{,} \PYG{l+m+mi}{761}\PYG{p}{,} \PYG{l+m+mi}{0}\PYG{p}{.}\PYG{l+m+mi}{1}\PYG{p}{,} \PYG{l+s+s1}{'length'}\PYG{p}{,} \PYG{k}{true}\PYG{p}{,} \PYG{k}{true}\PYG{p}{)}\PYG{p}{;}
\end{Verbatim}

\begin{Verbatim}[commandchars=\\\{\}]
  \PYG{n}{gid}   \PYG{o}{\textbar{}}                              \PYG{n}{the\PYGZus{}geom}
\PYG{c+c1}{--------+---------------------------------------------------------------}
    \PYG{l+m+mi}{293} \PYG{o}{\textbar{}} \PYG{n}{MULTILINESTRING}\PYG{p}{(}\PYG{p}{(}\PYG{l+m+mi}{18}\PYG{p}{.}\PYG{l+m+mi}{4074149} \PYG{o}{-}\PYG{l+m+mi}{33}\PYG{p}{.}\PYG{l+m+mi}{9443308}\PYG{p}{,}\PYG{l+m+mi}{18}\PYG{p}{.}\PYG{l+m+mi}{4074019} \PYG{o}{-}\PYG{l+m+mi}{33}\PYG{p}{.}\PYG{l+m+mi}{9443833}\PYG{p}{)}\PYG{p}{)}
    \PYG{l+m+mi}{293} \PYG{o}{\textbar{}} \PYG{n}{MULTILINESTRING}\PYG{p}{(}\PYG{p}{(}\PYG{l+m+mi}{18}\PYG{p}{.}\PYG{l+m+mi}{4074149} \PYG{o}{-}\PYG{l+m+mi}{33}\PYG{p}{.}\PYG{l+m+mi}{9443308}\PYG{p}{,}\PYG{l+m+mi}{18}\PYG{p}{.}\PYG{l+m+mi}{4074019} \PYG{o}{-}\PYG{l+m+mi}{33}\PYG{p}{.}\PYG{l+m+mi}{9443833}\PYG{p}{)}\PYG{p}{)}
   \PYG{l+m+mi}{4632} \PYG{o}{\textbar{}} \PYG{n}{MULTILINESTRING}\PYG{p}{(}\PYG{p}{(}\PYG{l+m+mi}{18}\PYG{p}{.}\PYG{l+m+mi}{4074149} \PYG{o}{-}\PYG{l+m+mi}{33}\PYG{p}{.}\PYG{l+m+mi}{9443308}\PYG{p}{,}\PYG{l+m+mi}{18}\PYG{p}{.}\PYG{l+m+mi}{4077388} \PYG{o}{-}\PYG{l+m+mi}{33}\PYG{p}{.}\PYG{l+m+mi}{9436183}\PYG{p}{)}\PYG{p}{)}
    \PYG{p}{.}\PYG{p}{.}\PYG{p}{.} \PYG{o}{\textbar{}} \PYG{p}{.}\PYG{p}{.}\PYG{p}{.}
    \PYG{l+m+mi}{762} \PYG{o}{\textbar{}} \PYG{n}{MULTILINESTRING}\PYG{p}{(}\PYG{p}{(}\PYG{l+m+mi}{18}\PYG{p}{.}\PYG{l+m+mi}{4241422} \PYG{o}{-}\PYG{l+m+mi}{33}\PYG{p}{.}\PYG{l+m+mi}{9179275}\PYG{p}{,}\PYG{l+m+mi}{18}\PYG{p}{.}\PYG{l+m+mi}{4237423} \PYG{o}{-}\PYG{l+m+mi}{33}\PYG{p}{.}\PYG{l+m+mi}{9182966}\PYG{p}{)}\PYG{p}{)}
    \PYG{l+m+mi}{761} \PYG{o}{\textbar{}} \PYG{n}{MULTILINESTRING}\PYG{p}{(}\PYG{p}{(}\PYG{l+m+mi}{18}\PYG{p}{.}\PYG{l+m+mi}{4243523} \PYG{o}{-}\PYG{l+m+mi}{33}\PYG{p}{.}\PYG{l+m+mi}{9177154}\PYG{p}{,}\PYG{l+m+mi}{18}\PYG{p}{.}\PYG{l+m+mi}{4241422} \PYG{o}{-}\PYG{l+m+mi}{33}\PYG{p}{.}\PYG{l+m+mi}{9179275}\PYG{p}{)}\PYG{p}{)}
\PYG{p}{(}\PYG{l+m+mi}{62} \PYG{k}{rows}\PYG{p}{)}
\end{Verbatim}

\begin{notice}{note}{Note:}
There is currently no wrapper function for A-Star without bounding box, since bounding boxes are very useful to increase performance. If you don't need a bounding box Dijkstra will be enough anyway.
\end{notice}

\begin{notice}{warning}{Warning:}
The projection of OSM data is ``degree'', so we set a bounding box containing start and end vertex plus a 0.1 degree buffer for example.
\end{notice}

\resetcurrentobjects
\hypertarget{--doc-chapters/advanced}{}

\chapter{Advanced usage of pgRouting shortest path search}

An ordinary shortest path query with result usualy looks like this:

\begin{Verbatim}[commandchars=\\\{\}]
\PYG{k}{SELECT} \PYG{o}{*} \PYG{k}{FROM} \PYG{n}{shortest\PYGZus{}path\PYGZus{}shooting\PYGZus{}star}\PYG{p}{(}
        \PYG{l+s+s1}{'SELECT gid as id, source, target, length as cost, x1, y1, x2, y2, rule,}
\PYG{l+s+s1}{        to\PYGZus{}cost, reverse\PYGZus{}cost FROM ways'}\PYG{p}{,} \PYG{l+m+mi}{1955}\PYG{p}{,} \PYG{l+m+mi}{5787}\PYG{p}{,} \PYG{k}{true}\PYG{p}{,} \PYG{k}{true}\PYG{p}{)}\PYG{p}{;}
\end{Verbatim}

Query result:

\begin{Verbatim}[commandchars=\\\{\}]
 \PYG{n}{vertex\PYGZus{}id} \PYG{o}{\textbar{}} \PYG{n}{edge\PYGZus{}id} \PYG{o}{\textbar{}}        \PYG{n}{cost}
\PYG{c+c1}{-----------+---------+---------------------}
      \PYG{l+m+mi}{8134} \PYG{o}{\textbar{}}    \PYG{l+m+mi}{1955} \PYG{o}{\textbar{}} \PYG{l+m+mi}{0}\PYG{p}{.}\PYG{l+m+mi}{00952475464810279}
      \PYG{l+m+mi}{5459} \PYG{o}{\textbar{}}    \PYG{l+m+mi}{1956} \PYG{o}{\textbar{}}  \PYG{l+m+mi}{0}\PYG{p}{.}\PYG{l+m+mi}{0628075563112871}
      \PYG{l+m+mi}{8137} \PYG{o}{\textbar{}}    \PYG{l+m+mi}{1976} \PYG{o}{\textbar{}}  \PYG{l+m+mi}{0}\PYG{p}{.}\PYG{l+m+mi}{0812786367080268}
      \PYG{l+m+mi}{5453} \PYG{o}{\textbar{}}     \PYG{l+m+mi}{758} \PYG{o}{\textbar{}}  \PYG{l+m+mi}{0}\PYG{p}{.}\PYG{l+m+mi}{0421747270358272}
      \PYG{l+m+mi}{5456} \PYG{o}{\textbar{}}    \PYG{l+m+mi}{3366} \PYG{o}{\textbar{}}  \PYG{l+m+mi}{0}\PYG{p}{.}\PYG{l+m+mi}{0104935732514831}
     \PYG{l+m+mi}{11086} \PYG{o}{\textbar{}}    \PYG{l+m+mi}{3367} \PYG{o}{\textbar{}}   \PYG{l+m+mi}{0}\PYG{p}{.}\PYG{l+m+mi}{113400030221047}
      \PYG{l+m+mi}{4416} \PYG{o}{\textbar{}}     \PYG{l+m+mi}{306} \PYG{o}{\textbar{}}   \PYG{l+m+mi}{0}\PYG{p}{.}\PYG{l+m+mi}{111600379959229}
      \PYG{l+m+mi}{4419} \PYG{o}{\textbar{}}     \PYG{l+m+mi}{307} \PYG{o}{\textbar{}}  \PYG{l+m+mi}{0}\PYG{p}{.}\PYG{l+m+mi}{0880411972519595}
      \PYG{l+m+mi}{4422} \PYG{o}{\textbar{}}    \PYG{l+m+mi}{4880} \PYG{o}{\textbar{}}  \PYG{l+m+mi}{0}\PYG{p}{.}\PYG{l+m+mi}{0208599114366633}
      \PYG{l+m+mi}{5101} \PYG{o}{\textbar{}}     \PYG{l+m+mi}{612} \PYG{o}{\textbar{}}  \PYG{l+m+mi}{0}\PYG{p}{.}\PYG{l+m+mi}{0906859882381495}
      \PYG{l+m+mi}{5102} \PYG{o}{\textbar{}}    \PYG{l+m+mi}{5787} \PYG{o}{\textbar{}}    \PYG{l+m+mi}{80089}\PYG{p}{.}\PYG{l+m+mi}{8820919459}
\PYG{p}{(}\PYG{l+m+mi}{11} \PYG{k}{rows}\PYG{p}{)}
\end{Verbatim}

That is usually called \emph{SHORTEST} path, which means that a length of an edge is its cost.

\textbf{Costs can be anything} (``Weighted costs'')

But in real networks we have different limitations or preferences for different road types for example. In other words, we want to calculate CHEAPEST path - a path with a minimal cost. There is no limitation in what we take as costs.

When we convert data from OSM format using the osm2pgrouting tool, we get these two additional tables for road types and classes:

\begin{Verbatim}[commandchars=@\[\]]
\d classes

  id @textbar[]   name
-----+------------
   2 @textbar[] cycleway
   1 @textbar[] highway
   4 @textbar[] junction
   3 @textbar[] tracktype
\end{Verbatim}

\begin{Verbatim}[commandchars=@\[\]]
\d types

 id  @textbar[] type@_id @textbar[]        name        @textbar[]  cost
-----+---------+--------------------+--------
 201 @textbar[]       2 @textbar[] lane               @textbar[]   1
 204 @textbar[]       2 @textbar[] opposite           @textbar[]   1
 203 @textbar[]       2 @textbar[] opposite@_lane      @textbar[]   1
 202 @textbar[]       2 @textbar[] track              @textbar[]   1
 117 @textbar[]       1 @textbar[] bridleway          @textbar[]   1
 113 @textbar[]       1 @textbar[] bus@_guideway       @textbar[]   1
 118 @textbar[]       1 @textbar[] byway              @textbar[]   1
 115 @textbar[]       1 @textbar[] cicleway           @textbar[]   1
 116 @textbar[]       1 @textbar[] footway            @textbar[]   1
 108 @textbar[]       1 @textbar[] living@_street      @textbar[]   1
 101 @textbar[]       1 @textbar[] motorway           @textbar[]   0.2
 103 @textbar[]       1 @textbar[] motorway@_junction  @textbar[]   0.2
 102 @textbar[]       1 @textbar[] motorway@_link      @textbar[]   0.2
 114 @textbar[]       1 @textbar[] path               @textbar[]   100
 111 @textbar[]       1 @textbar[] pedestrian         @textbar[]   100
 106 @textbar[]       1 @textbar[] primary            @textbar[]   100
 107 @textbar[]       1 @textbar[] primary@_link       @textbar[]   100
 107 @textbar[]       1 @textbar[] residential        @textbar[]   100
 100 @textbar[]       1 @textbar[] road               @textbar[]   0.7
 100 @textbar[]       1 @textbar[] unclassified       @textbar[]   0.7
 106 @textbar[]       1 @textbar[] secondary          @textbar[]   10
 109 @textbar[]       1 @textbar[] service            @textbar[]   10
 112 @textbar[]       1 @textbar[] services           @textbar[]   10
 119 @textbar[]       1 @textbar[] steps              @textbar[]   10
 107 @textbar[]       1 @textbar[] tertiary           @textbar[]   10
 110 @textbar[]       1 @textbar[] track              @textbar[]   10
 104 @textbar[]       1 @textbar[] trunk              @textbar[]   10
 105 @textbar[]       1 @textbar[] trunk@_link         @textbar[]   10
 401 @textbar[]       4 @textbar[] roundabout         @textbar[]   10
 301 @textbar[]       3 @textbar[] grade1             @textbar[]   15
 302 @textbar[]       3 @textbar[] grade2             @textbar[]   15
 303 @textbar[]       3 @textbar[] grade3             @textbar[]   15
 304 @textbar[]       3 @textbar[] grade4             @textbar[]   15
 305 @textbar[]       3 @textbar[] grade5             @textbar[]   15
\end{Verbatim}

Road class is linked with the ways table by class\_id field. Cost values for classes table are assigned arbitrary.

\begin{Verbatim}[commandchars=\\\{\}]
\PYG{k}{UPDATE} \PYG{n}{classes} \PYG{k}{SET} \PYG{n}{cost}\PYG{o}{=}\PYG{l+m+mi}{15} \PYG{k}{WHERE} \PYG{n}{id}\PYG{o}{\textgreater{}}\PYG{l+m+mi}{300}\PYG{p}{;}
\end{Verbatim}

For better performance it is worth to create an index on id field of classes table.

\begin{Verbatim}[commandchars=\\\{\}]
\PYG{k}{CREATE} \PYG{k}{INDEX} \PYG{n}{class\PYGZus{}idx} \PYG{k}{ON} \PYG{n}{ways} \PYG{p}{(}\PYG{n}{id}\PYG{p}{)}\PYG{p}{;}
\end{Verbatim}

The idea behind these two tables is to specify a factor to be multiplied with the cost of each link (usually length):

\begin{Verbatim}[commandchars=\\\{\}]
\PYG{k}{SELECT} \PYG{o}{*} \PYG{k}{FROM} \PYG{n}{shortest\PYGZus{}path\PYGZus{}shooting\PYGZus{}star}\PYG{p}{(}
        \PYG{l+s+s1}{'SELECT gid as id, class\PYGZus{}id, source, target, length*c.cost as cost,}
\PYG{l+s+s1}{                x1, y1, x2, y2, rule, to\PYGZus{}cost, reverse\PYGZus{}cost*c.cost as reverse\PYGZus{}cost}
\PYG{l+s+s1}{        FROM ways w, classes c}
\PYG{l+s+s1}{        WHERE class\PYGZus{}id=c.id'}\PYG{p}{,} \PYG{l+m+mi}{1955}\PYG{p}{,} \PYG{l+m+mi}{5787}\PYG{p}{,} \PYG{k}{true}\PYG{p}{,} \PYG{k}{true}\PYG{p}{)}\PYG{p}{;}
\end{Verbatim}

Query result:

\begin{Verbatim}[commandchars=\\\{\}]
 \PYG{n}{vertex\PYGZus{}id} \PYG{o}{\textbar{}} \PYG{n}{edge\PYGZus{}id} \PYG{o}{\textbar{}}        \PYG{n}{cost}
\PYG{c+c1}{-----------+---------+---------------------}
      \PYG{l+m+mi}{8134} \PYG{o}{\textbar{}}    \PYG{l+m+mi}{1955} \PYG{o}{\textbar{}} \PYG{l+m+mi}{0}\PYG{p}{.}\PYG{l+m+mi}{00666732825367195}
      \PYG{l+m+mi}{5459} \PYG{o}{\textbar{}}    \PYG{l+m+mi}{1956} \PYG{o}{\textbar{}}   \PYG{l+m+mi}{0}\PYG{p}{.}\PYG{l+m+mi}{043965289417901}
      \PYG{l+m+mi}{8137} \PYG{o}{\textbar{}}    \PYG{l+m+mi}{1992} \PYG{o}{\textbar{}}   \PYG{l+m+mi}{0}\PYG{p}{.}\PYG{l+m+mi}{126646230936747}
      \PYG{l+m+mi}{5464} \PYG{o}{\textbar{}}     \PYG{l+m+mi}{762} \PYG{o}{\textbar{}}   \PYG{l+m+mi}{0}\PYG{p}{.}\PYG{l+m+mi}{827868704808978}
      \PYG{l+m+mi}{5467} \PYG{o}{\textbar{}}     \PYG{l+m+mi}{763} \PYG{o}{\textbar{}}    \PYG{l+m+mi}{0}\PYG{p}{.}\PYG{l+m+mi}{16765902528648}
      \PYG{p}{.}\PYG{p}{.}\PYG{p}{.} \PYG{o}{\textbar{}}     \PYG{p}{.}\PYG{p}{.}\PYG{p}{.} \PYG{o}{\textbar{}}                 \PYG{p}{.}\PYG{p}{.}\PYG{p}{.}
      \PYG{l+m+mi}{9790} \PYG{o}{\textbar{}}    \PYG{l+m+mi}{5785} \PYG{o}{\textbar{}} \PYG{l+m+mi}{0}\PYG{p}{.}\PYG{l+m+mi}{00142107468268373}
      \PYG{l+m+mi}{8548} \PYG{o}{\textbar{}}    \PYG{l+m+mi}{5786} \PYG{o}{\textbar{}} \PYG{l+m+mi}{0}\PYG{p}{.}\PYG{l+m+mi}{00066608685984761}
      \PYG{l+m+mi}{16214} \PYG{o}{\textbar{}}    \PYG{l+m+mi}{5787} \PYG{o}{\textbar{}}  \PYG{l+m+mi}{0}\PYG{p}{.}\PYG{l+m+mi}{0160179764183892}
\PYG{p}{(}\PYG{l+m+mi}{69} \PYG{k}{rows}\PYG{p}{)}
\end{Verbatim}

We can see that the shortest path result is completely different from the example before. We call this ``weighted costs''.

Another example is to restrict access to roads of a certain type:

\begin{Verbatim}[commandchars=\\\{\}]
\PYG{k}{UPDATE} \PYG{n}{classes} \PYG{k}{SET} \PYG{n}{cost}\PYG{o}{=}\PYG{l+m+mi}{100000} \PYG{k}{WHERE} \PYG{n}{name} \PYG{k}{LIKE} \PYG{l+s+s1}{'motorway\%'}\PYG{p}{;}
\end{Verbatim}

Through subqueries you can ``mix'' your costs as you like and this will change the results of your routing request immediately. Cost changes will affect the next shortest path search, and there is no need to rebuild your network.

\resetcurrentobjects
\hypertarget{--doc-chapters/php_server}{}

\chapter{Server side scripts with PHP}

We will use a PHP script to make the routing query and send the result back to the web client.

The following steps are necessary:

Retrieve the start and end point coordinates.
Find the closest edge to start/end point.
Take either the start or end vertex of this edge (for Dijkstra/ A-Star) or the complete edge (Shooting-Star) as start of the route and end respectively.
Make the Shortest Path database query.
Transform the query result to XML and send it back to the web client.

\resetcurrentobjects
\hypertarget{--doc-chapters/geoext_client}{}

\chapter{GeoExt Browser Client}


\renewcommand{\indexname}{Module Index}

\renewcommand{\indexname}{Index}
\printindex
\end{document}
